% p075_2; source PDF page 75; manually verified editable table.
\begin{tikzpicture}[x=1cm,y=-1cm,inner sep=0pt,outer sep=0pt]
\path[use as bounding box] (0,0) rectangle (16.000,5.985);
\draw[line width=.4pt] (0.000,0.000) rectangle (4.300,0.855);
\node[align=center,text width=4.040cm,font={\linespread{1}\fontsize{10}{13.50}\selectfont\bfseries}] at (2.150,0.427) {事项};
\draw[line width=.4pt] (4.300,0.000) rectangle (10.100,0.855);
\node[align=center,text width=5.540cm,font={\linespread{1}\fontsize{10}{13.50}\selectfont\bfseries}] at (7.200,0.427) {故意犯罪};
\draw[line width=.4pt] (10.100,0.000) rectangle (16.000,0.855);
\node[align=center,text width=5.640cm,font={\linespread{1}\fontsize{10}{13.50}\selectfont\bfseries}] at (13.050,0.427) {过失犯罪};
\draw[line width=.4pt] (0.000,0.855) rectangle (4.300,1.710);
\node[align=left,text width=4.040cm,font={\linespread{1}\fontsize{10}{13.50}\selectfont}] at (2.150,1.282) {成立条件};
\draw[line width=.4pt] (4.300,0.855) rectangle (10.100,1.710);
\node[align=left,text width=5.540cm,font={\linespread{1}\fontsize{10}{13.50}\selectfont}] at (7.200,1.282) {一般不要求发生实害结果};
\draw[line width=.4pt] (10.100,0.855) rectangle (16.000,1.710);
\node[align=left,text width=5.640cm,font={\linespread{1}\fontsize{10}{13.50}\selectfont}] at (13.050,1.282) {要求发生实害结果};
\draw[line width=.4pt] (0.000,1.710) rectangle (4.300,2.565);
\node[align=left,text width=4.040cm,font={\linespread{1}\fontsize{10}{13.50}\selectfont}] at (2.150,2.138) {主客观相一致原则};
\draw[line width=.4pt] (4.300,1.710) rectangle (10.100,2.565);
\node[align=left,text width=5.540cm,font={\linespread{1}\fontsize{10}{13.50}\selectfont}] at (7.200,2.138) {要求遵守主客观相一致原则};
\draw[line width=.4pt] (10.100,1.710) rectangle (16.000,2.565);
\node[align=left,text width=5.640cm,font={\linespread{1}\fontsize{10}{13.50}\selectfont}] at (13.050,2.138) {不存在遵守该原则的必要};
\draw[line width=.4pt] (0.000,2.565) rectangle (4.300,3.420);
\node[align=left,text width=4.040cm,font={\linespread{1}\fontsize{10}{13.50}\selectfont}] at (2.150,2.992) {主观恶性程度};
\draw[line width=.4pt] (4.300,2.565) rectangle (10.100,3.420);
\node[align=left,text width=5.540cm,font={\linespread{1}\fontsize{10}{13.50}\selectfont}] at (7.200,2.992) {较重，因此刑罚也较重};
\draw[line width=.4pt] (10.100,2.565) rectangle (16.000,3.420);
\node[align=left,text width=5.640cm,font={\linespread{1}\fontsize{10}{13.50}\selectfont}] at (13.050,2.992) {较轻，因此刑罚也较轻};
\draw[line width=.4pt] (0.000,3.420) rectangle (4.300,4.275);
\node[align=left,text width=4.040cm,font={\linespread{1}\fontsize{10}{13.50}\selectfont}] at (2.150,3.848) {未完成形态};
\draw[line width=.4pt] (4.300,3.420) rectangle (10.100,4.275);
\node[align=left,text width=5.540cm,font={\linespread{1}\fontsize{10}{13.50}\selectfont}] at (7.200,3.848) {有（包括间接故意）};
\draw[line width=.4pt] (10.100,3.420) rectangle (16.000,4.275);
\node[align=left,text width=5.640cm,font={\linespread{1}\fontsize{10}{13.50}\selectfont}] at (13.050,3.848) {无，只有成立与否问题};
\draw[line width=.4pt] (0.000,4.275) rectangle (4.300,5.130);
\node[align=left,text width=4.040cm,font={\linespread{1}\fontsize{10}{13.50}\selectfont}] at (2.150,4.703) {共同犯罪};
\draw[line width=.4pt] (4.300,4.275) rectangle (10.100,5.130);
\node[align=left,text width=5.540cm,font={\linespread{1}\fontsize{10}{13.50}\selectfont}] at (7.200,4.703) {有};
\draw[line width=.4pt] (10.100,4.275) rectangle (16.000,5.130);
\node[align=left,text width=5.640cm,font={\linespread{1}\fontsize{10}{13.50}\selectfont}] at (13.050,4.703) {无};
\draw[line width=.4pt] (0.000,5.130) rectangle (4.300,5.985);
\node[align=left,text width=4.040cm,font={\linespread{1}\fontsize{10}{13.50}\selectfont}] at (2.150,5.558) {处罚};
\draw[line width=.4pt] (4.300,5.130) rectangle (10.100,5.985);
\node[align=left,text width=5.540cm,font={\linespread{1}\fontsize{10}{13.50}\selectfont}] at (7.200,5.558) {处罚故意犯罪是原则};
\draw[line width=.4pt] (10.100,5.130) rectangle (16.000,5.985);
\node[align=left,text width=5.640cm,font={\linespread{1}\fontsize{10}{13.50}\selectfont}] at (13.050,5.558) {处罚过失犯罪是例外};
\end{tikzpicture}
